\documentclass[12pt]{article}
\usepackage[margin=1in]{geometry}
\usepackage{amsmath,amssymb,amsthm}
\usepackage{bm}
\usepackage{bbold}
\usepackage{booktabs}
\usepackage{enumitem}
\usepackage{tcolorbox}
\usepackage{fancyhdr}
\usepackage{hyperref}

\pagestyle{fancy}
\lhead{PLSC 30700: Estimation I}
\rhead{Robert Gulotty}
\chead{}
\cfoot{\thepage}

\newcommand{\E}{\mathbb{E}}
\newcommand{\Var}{\operatorname{Var}}
\newcommand{\Cov}{\operatorname{Cov}}
\newcommand{\bfX}{\mathbf{X}}
\newcommand{\bfx}{\mathbf{x}}
\newcommand{\bfy}{\mathbf{y}}
\newcommand{\bfbeta}{\bm{\beta}}
\newcommand{\bfQ}{\mathbf{Q}}

\theoremstyle{definition}
\newtheorem{step}{Step}
\newtheorem{remark}{Remark}

\title{\textbf{Linear Algebra Walkthrough}\\[4pt]
\large Lecture 2: CEF and Best Linear Predictor}
\author{Companion to Slides --- PLSC 30700}
\date{}

\begin{document}
\maketitle
\thispagestyle{fancy}

\noindent This document expands every linear algebra step in Lecture 2 (``Linear Models: CEF and Best Linear Predictor''), connecting each derivation back to the corresponding slide. Where the slides use compact matrix notation, we unpack the matrices into their column-vector and element-wise forms.

\tableofcontents
\newpage

%======================================================================
\section{Bivariate BLP: Solving the First-Order Conditions}
\label{sec:bivariate}
%======================================================================

\begin{tcolorbox}[colback=blue!5, colframe=blue!50, title=Slide Reference]
``Solving for Best Linear Predictor ($a$ and $b$)''
\end{tcolorbox}

We want to find scalars $a$ and $b$ that minimize the mean squared prediction error:
\[
M(a,b) = \E\!\left[(Y - (a + bX))^2\right].
\]

\subsection{Expanding the square}

Write out the square explicitly:
\begin{align}
M(a,b) &= \E\!\left[Y^2 - 2Y(a+bX) + (a+bX)^2\right] \nonumber\\
&= \E[Y^2] - 2a\E[Y] - 2b\E[XY] + a^2 + 2ab\E[X] + b^2\E[X^2]. \label{eq:Mab}
\end{align}
Every term uses the linearity of expectation: $\E[cZ] = c\E[Z]$ for a constant $c$.

\subsection{FOC with respect to $a$}

Differentiate \eqref{eq:Mab} with respect to $a$, holding $b$ fixed:
\[
\frac{\partial M}{\partial a} = -2\E[Y] + 2a + 2b\E[X] = 0.
\]
Solving:
\begin{equation}
\boxed{a = \E[Y] - b\E[X].} \label{eq:astar}
\end{equation}
\textbf{Interpretation.}  The intercept adjusts the prediction so that, on average, $a + b\E[X] = \E[Y]$---the predictor passes through the point $(\E[X], \E[Y])$.

\subsection{FOC with respect to $b$}

Differentiate \eqref{eq:Mab} with respect to $b$:
\[
\frac{\partial M}{\partial b} = -2\E[XY] + 2a\E[X] + 2b\E[X^2] = 0.
\]
Substitute $a = \E[Y] - b\E[X]$ from \eqref{eq:astar}:
\begin{align*}
\E[XY] &= \bigl(\E[Y] - b\E[X]\bigr)\E[X] + b\E[X^2]\\
&= \E[Y]\E[X] - b(\E[X])^2 + b\E[X^2]\\
&= \E[Y]\E[X] + b\bigl(\E[X^2] - (\E[X])^2\bigr).
\end{align*}
Rearranging:
\[
\E[XY] - \E[X]\E[Y] = b\bigl(\E[X^2] - (\E[X])^2\bigr).
\]
The left side is $\Cov(X,Y)$ and the right side involves $\Var(X)$:
\begin{equation}
\boxed{b = \frac{\Cov(X,Y)}{\Var(X)}.} \label{eq:bstar}
\end{equation}

\subsection{The BLP in one formula}

Substituting back:
\begin{align*}
\E^*[Y|X] &= a + bX = \E[Y] - b\E[X] + bX\\
&= \E[Y] + \frac{\Cov(X,Y)}{\Var(X)}\bigl(X - \E[X]\bigr).
\end{align*}
This is the population regression line, centered at $(\E[X], \E[Y])$ with slope $\Cov(X,Y)/\Var(X)$.


%======================================================================
\section{Variance of the BLP Error (Bivariate)}
\label{sec:bivar-var}
%======================================================================

\begin{tcolorbox}[colback=blue!5, colframe=blue!50, title=Slide Reference]
``Deriving Variance of BLP at Minimized Values''
\end{tcolorbox}

Define the projection error $e = Y - (\alpha + \beta X)$, where $\alpha = \E[Y] - \beta\E[X]$. We want $\Var(e)$.

\begin{step}[Rewrite the error]
Since $\alpha = \E[Y] - \beta\E[X]$, subtracting $\alpha + \beta X$ from $Y$ gives:
\[
e = Y - \E[Y] + \beta\E[X] - \beta X = (Y - \E[Y]) - \beta(X - \E[X]).
\]
So the error is a difference of two centered random variables.
\end{step}

\begin{step}[Compute the variance]
Using $\Var(A - B) = \Var(A) + \Var(B) - 2\Cov(A,B)$ with $A = Y$ and $B = \beta X$:
\begin{align*}
\Var(e) &= \Var(Y - \beta X) \\
&= \Var(Y) + \beta^2\Var(X) - 2\beta\Cov(X,Y).
\end{align*}
The constant $\alpha$ drops out because $\Var(Z + c) = \Var(Z)$.
\end{step}

\begin{step}[Substitute $\beta\Var(X) = \Cov(X,Y)$]
From equation~\eqref{eq:bstar}, $\beta = \Cov(X,Y)/\Var(X)$, so $\beta\Var(X) = \Cov(X,Y)$. Thus:
\begin{align*}
\Var(e) &= \Var(Y) + \beta^2\Var(X) - 2\beta \cdot \beta\Var(X)\\
&= \Var(Y) + \beta^2\Var(X) - 2\beta^2\Var(X)\\
&= \Var(Y) - \beta^2\Var(X).
\end{align*}
\end{step}

\begin{equation}
\boxed{\sigma_e^2 = \sigma_Y^2 - \beta^2\sigma_X^2.}
\end{equation}

\textbf{Interpretation.}  The error variance is the total variance of $Y$ minus what the linear predictor explains.  The quantity $\beta^2\sigma_X^2 / \sigma_Y^2$ is the population $R^2$.

%======================================================================
\section{Multivariate BLP: Expanding the MSPE}
\label{sec:multivariate}
%======================================================================

\begin{tcolorbox}[colback=blue!5, colframe=blue!50, title=Slide Reference]
``Deriving the Linear Projection Coefficient $\beta$''
\end{tcolorbox}

Now $X$ is a $k \times 1$ random vector and $\beta$ is a $k \times 1$ parameter vector. To make the algebra concrete, write:
\[
X = \begin{pmatrix} X_1 \\ X_2 \\ \vdots \\ X_{k-1} \\ 1 \end{pmatrix}, \qquad
\beta = \begin{pmatrix} \beta_1 \\ \beta_2 \\ \vdots \\ \beta_{k-1} \\ \beta_k \end{pmatrix}.
\]
The last element of $X$ is the constant 1 (for the intercept), so $\beta_k$ plays the role of the intercept.

The linear predictor is the inner product:
\[
X'\beta = \begin{pmatrix} X_1 & X_2 & \cdots & X_{k-1} & 1 \end{pmatrix}
\begin{pmatrix} \beta_1 \\ \beta_2 \\ \vdots \\ \beta_{k-1} \\ \beta_k \end{pmatrix}
= \beta_1 X_1 + \beta_2 X_2 + \cdots + \beta_{k-1} X_{k-1} + \beta_k.
\]

\subsection{Step 1: Writing the MSPE}

\begin{align}
S(\beta) &= \E\!\left[(Y - X'\beta)^2\right] = \E\!\left[(Y - X'\beta)(Y - X'\beta)\right]. \label{eq:Sbeta}
\end{align}

\subsection{Step 2: Expanding the product}

Note that $Y - X'\beta$ is a \emph{scalar} random variable ($1 \times 1$), so $(Y - X'\beta)^2 = (Y - X'\beta)(Y - X'\beta)$. Expand:
\begin{align*}
(Y - X'\beta)^2 &= Y^2 - 2Y \cdot X'\beta + (X'\beta)^2.
\end{align*}

We need to understand each term:

\paragraph{Term 1: $Y^2$.} Just the square of the scalar $Y$.

\paragraph{Term 2: $Y \cdot X'\beta$.} Since $X'\beta$ is a scalar:
\[
Y \cdot X'\beta = Y(\beta_1 X_1 + \beta_2 X_2 + \cdots + \beta_k) = \beta_1 YX_1 + \beta_2 YX_2 + \cdots + \beta_k Y.
\]
Taking expectations and writing in matrix form:
\[
\E[Y \cdot X'\beta] = \beta' \E[XY],
\]
where $\E[XY]$ is the $k \times 1$ vector:
\[
\E[XY] = \begin{pmatrix} \E[X_1 Y] \\ \E[X_2 Y] \\ \vdots \\ \E[Y] \end{pmatrix}.
\]
(The last entry is $\E[1 \cdot Y] = \E[Y]$.)

\paragraph{Term 3: $(X'\beta)^2$.} Since $X'\beta$ is a scalar, $(X'\beta)^2 = (X'\beta)(X'\beta) = (\beta'X)(X'\beta) = \beta'(XX')\beta$.

The matrix $XX'$ is $k \times k$. Expanding it column-by-column:
\[
XX' = \begin{pmatrix} X_1 \\ X_2 \\ \vdots \\ 1 \end{pmatrix}
\begin{pmatrix} X_1 & X_2 & \cdots & 1 \end{pmatrix}
= \begin{pmatrix}
X_1^2 & X_1 X_2 & \cdots & X_1 \\
X_2 X_1 & X_2^2 & \cdots & X_2 \\
\vdots & \vdots & \ddots & \vdots \\
X_1 & X_2 & \cdots & 1
\end{pmatrix}.
\]

Taking expectations element-by-element:
\[
\E[XX'] = \bfQ_{XX} = \begin{pmatrix}
\E[X_1^2] & \E[X_1 X_2] & \cdots & \E[X_1] \\
\E[X_2 X_1] & \E[X_2^2] & \cdots & \E[X_2] \\
\vdots & \vdots & \ddots & \vdots \\
\E[X_1] & \E[X_2] & \cdots & 1
\end{pmatrix}.
\]

\textbf{Key fact:} $\bfQ_{XX}$ is symmetric ($\E[X_i X_j] = \E[X_j X_i]$) and, by the regularity condition, positive definite.

\subsection{Step 3: The MSPE in matrix form}

Putting the three terms together under the expectation:
\begin{equation}
S(\beta) = \E[Y^2] - 2\beta'\E[XY] + \beta'\E[XX']\beta. \label{eq:Smatrix}
\end{equation}

\begin{remark}
Notice the parallel with the bivariate case in Section~\ref{sec:bivariate}. The scalar $\E[X^2]$ has become the matrix $\E[XX']$, and the scalar $\E[XY]$ has become the vector $\E[XY]$.
\end{remark}

\subsection{Step 4: Differentiating with respect to $\beta$}

We use the matrix calculus rules:
\begin{itemize}
\item $\frac{\partial}{\partial \beta}(\beta' \mathbf{a}) = \mathbf{a}$ \quad for constant vector $\mathbf{a}$.
\item $\frac{\partial}{\partial \beta}(\beta' \mathbf{A} \beta) = 2\mathbf{A}\beta$ \quad for symmetric matrix $\mathbf{A}$.
\end{itemize}

Applying these to \eqref{eq:Smatrix}:
\begin{align*}
\frac{\partial S}{\partial \beta} &= -2\E[XY] + 2\E[XX']\beta.
\end{align*}

\textbf{Element by element.}  The $j$-th component of the gradient is:
\[
\frac{\partial S}{\partial \beta_j} = -2\E[X_j Y] + 2\sum_{i=1}^{k} \E[X_j X_i]\beta_i, \quad j = 1, \ldots, k.
\]
This is a system of $k$ equations in $k$ unknowns.

\subsection{Step 5: Setting the FOC to zero and solving}

\begin{align}
0 &= -2\E[XY] + 2\E[XX']\beta \nonumber\\
\E[XY] &= \E[XX']\beta \nonumber\\
\bfQ_{XY} &= \bfQ_{XX}\beta. \label{eq:normalQ}
\end{align}

Since $\bfQ_{XX}$ is positive definite (regularity condition 3), it is invertible. Premultiply both sides by $\bfQ_{XX}^{-1}$:
\begin{equation}
\boxed{\beta = \bfQ_{XX}^{-1}\bfQ_{XY} = \E[XX']^{-1}\E[XY].} \label{eq:betastar}
\end{equation}

\textbf{Written out for $k=3$ (two regressors plus intercept):}
\[
\begin{pmatrix} \beta_1 \\ \beta_2 \\ \beta_3 \end{pmatrix}
=
\begin{pmatrix}
\E[X_1^2] & \E[X_1 X_2] & \E[X_1] \\
\E[X_1 X_2] & \E[X_2^2] & \E[X_2] \\
\E[X_1] & \E[X_2] & 1
\end{pmatrix}^{-1}
\begin{pmatrix}
\E[X_1 Y] \\
\E[X_2 Y] \\
\E[Y]
\end{pmatrix}.
\]

%======================================================================
\section{The Projection Error is Uncorrelated with $X$}
\label{sec:orthogonality}
%======================================================================

\begin{tcolorbox}[colback=blue!5, colframe=blue!50, title=Slide Reference]
``Linear Projection and Projection Error''
\end{tcolorbox}

Define $e = Y - X'\beta$, so $Y = X'\beta + e$. We show $\E[Xe] = \mathbf{0}$.

\subsection{The algebra}

\begin{align*}
\E[Xe] &= \E[X(Y - X'\beta)]\\
&= \E[XY] - \E[XX']\beta.
\end{align*}

\textbf{Why $\E[X \cdot X'\beta] = \E[XX']\beta$:} The product $X \cdot X'\beta$ has $j$-th element:
\[
X_j \cdot (X'\beta) = X_j \sum_{i=1}^k X_i \beta_i = \sum_{i=1}^k X_j X_i \beta_i.
\]
Taking expectations: $\E[X_j \cdot X'\beta] = \sum_i \E[X_j X_i]\beta_i$, which is the $j$-th row of $\E[XX']\beta$.

Now substitute $\beta = \E[XX']^{-1}\E[XY]$:
\begin{align*}
\E[Xe] &= \E[XY] - \E[XX'] \cdot \E[XX']^{-1}\E[XY]\\
&= \E[XY] - \mathbf{I}_k \cdot \E[XY]\\
&= \E[XY] - \E[XY] = \mathbf{0}.
\end{align*}

\subsection{What these $k$ equations say}

Written as a vector equation:
\[
\E[Xe] = \begin{pmatrix} \E[X_1 e] \\ \E[X_2 e] \\ \vdots \\ \E[1 \cdot e] \end{pmatrix} = \begin{pmatrix} 0 \\ 0 \\ \vdots \\ 0 \end{pmatrix}.
\]

\begin{itemize}
\item \textbf{Equations 1 through $k{-}1$:} $\E[X_j e] = 0$ for each regressor. The error is uncorrelated with every regressor.
\item \textbf{Equation $k$ (the constant):} $\E[1 \cdot e] = \E[e] = 0$. The error has mean zero.
\end{itemize}

This is why including a constant matters: it forces $\E[e] = 0$.


%======================================================================
\section{Why $\bfQ_{XX}$ Must Be Positive Definite}
\label{sec:pd}
%======================================================================

\begin{tcolorbox}[colback=blue!5, colframe=blue!50, title=Slide Reference]
``The Design Matrix''
\end{tcolorbox}

The BLP requires inverting $\bfQ_{XX} = \E[XX']$. We need this matrix to be positive definite.

\subsection{Positive semi-definiteness is automatic}

For any non-zero vector $\alpha \in \mathbb{R}^k$:
\[
\alpha'\bfQ_{XX}\alpha = \alpha'\E[XX']\alpha = \E[\alpha'XX'\alpha] = \E[(\alpha'X)^2] \geq 0.
\]

\textbf{Expanding $\alpha'X$:}
\[
\alpha'X = \alpha_1 X_1 + \alpha_2 X_2 + \cdots + \alpha_k \cdot 1.
\]
This is a linear combination of the regressors. Its square is non-negative, so its expectation is non-negative.

\subsection{Strict positive definiteness = no perfect collinearity}

$\bfQ_{XX}$ is positive definite (not just semi-definite) if and only if $\E[(\alpha'X)^2] > 0$ for all $\alpha \neq \mathbf{0}$.

If $\E[(\alpha'X)^2] = 0$ for some $\alpha \neq \mathbf{0}$, then $\alpha'X = 0$ almost surely, meaning:
\[
\alpha_1 X_1 + \alpha_2 X_2 + \cdots + \alpha_k = 0 \quad \text{with probability 1}.
\]
This is \textbf{perfect collinearity}: one regressor is an exact linear function of the others. In that case $\bfQ_{XX}$ is singular, $\beta$ is not uniquely defined, and the BLP does not exist.


%======================================================================
\section{Multivariate Error Variance}
\label{sec:errorvar}
%======================================================================

\begin{tcolorbox}[colback=blue!5, colframe=blue!50, title=Slide Reference]
``Linear Predictor Error Variance''
\end{tcolorbox}

Define $\sigma^2 = \E[e^2]$ where $e = Y - X'\beta$. We expand this using the matrices $Q_{YY} = \E[Y^2]$, $\bfQ_{YX} = \E[YX']$ (a $1 \times k$ row vector), and $\bfQ_{XY} = \E[XY]$ (a $k \times 1$ column vector).

\subsection{Expanding $\E[e^2]$}

\begin{align*}
\sigma^2 &= \E[(Y - X'\beta)^2]\\
&= \E[Y^2] - 2\E[YX']\beta + \beta'\E[XX']\beta\\
&= Q_{YY} - 2\bfQ_{YX}\beta + \beta'\bfQ_{XX}\beta.
\end{align*}

\subsection{Substituting $\beta = \bfQ_{XX}^{-1}\bfQ_{XY}$}

\paragraph{The second term:}
\[
2\bfQ_{YX}\beta = 2\bfQ_{YX}\bfQ_{XX}^{-1}\bfQ_{XY}.
\]
(This is a $1 \times k$ times $k \times k$ times $k \times 1 = $ scalar.)

\paragraph{The third term:}
\begin{align*}
\beta'\bfQ_{XX}\beta &= (\bfQ_{XX}^{-1}\bfQ_{XY})'\bfQ_{XX}(\bfQ_{XX}^{-1}\bfQ_{XY})\\
&= \bfQ_{YX}(\bfQ_{XX}^{-1})'\bfQ_{XX}\bfQ_{XX}^{-1}\bfQ_{XY}.
\end{align*}
Since $\bfQ_{XX}$ is symmetric, $(\bfQ_{XX}^{-1})' = \bfQ_{XX}^{-1}$, so:
\[
\beta'\bfQ_{XX}\beta = \bfQ_{YX}\bfQ_{XX}^{-1}\underbrace{\bfQ_{XX}\bfQ_{XX}^{-1}}_{\mathbf{I}_k}\bfQ_{XY} = \bfQ_{YX}\bfQ_{XX}^{-1}\bfQ_{XY}.
\]

\paragraph{Combining:}
\begin{align*}
\sigma^2 &= Q_{YY} - 2\bfQ_{YX}\bfQ_{XX}^{-1}\bfQ_{XY} + \bfQ_{YX}\bfQ_{XX}^{-1}\bfQ_{XY}\\
&= Q_{YY} - \bfQ_{YX}\bfQ_{XX}^{-1}\bfQ_{XY}.
\end{align*}

\begin{equation}
\boxed{\sigma^2 \equiv Q_{YY \cdot X} = Q_{YY} - \bfQ_{YX}\bfQ_{XX}^{-1}\bfQ_{XY}.}
\end{equation}

The notation $Q_{YY\cdot X}$ is read ``the variance of $Y$ after projecting out $X$.''

\subsection{Expanded for $k=2$ (one regressor plus intercept)}

Let $X = (X_1, 1)'$. Then:
\[
\bfQ_{YX}\bfQ_{XX}^{-1}\bfQ_{XY} =
\begin{pmatrix} \E[X_1 Y] & \E[Y] \end{pmatrix}
\begin{pmatrix} \E[X_1^2] & \E[X_1] \\ \E[X_1] & 1 \end{pmatrix}^{-1}
\begin{pmatrix} \E[X_1 Y] \\ \E[Y] \end{pmatrix}.
\]
This is the multivariate generalization of $\beta^2\Var(X)$ from Section~\ref{sec:bivar-var}.


%======================================================================
\section{Intercepts, Slopes, and Demeaning}
\label{sec:demeaning}
%======================================================================

\begin{tcolorbox}[colback=blue!5, colframe=blue!50, title=Slide Reference]
``Intercepts and Slopes''
\end{tcolorbox}

The slides separate the constant from $X$ and write $Y = X'\beta + \alpha + e$, then show that demeaning produces $\beta = \Var[X]^{-1}\Cov(X,Y)$.

\subsection{Setup: separating the intercept}

Let $X = (X_1, \ldots, X_{k-1})'$ (no constant) and write:
\[
Y = X'\beta + \alpha + e.
\]
From the FOC for $\alpha$ (analogous to Section~\ref{sec:bivariate}):
\[
\alpha = \E[Y] - \E[X]'\beta = \mu_Y - \mu_X'\beta.
\]

\subsection{Demeaning}

Substitute $\alpha$:
\begin{align*}
Y &= X'\beta + (\mu_Y - \mu_X'\beta) + e\\
Y - \mu_Y &= (X - \mu_X)'\beta + e.
\end{align*}

Now define the centered variables $\tilde{Y} = Y - \mu_Y$ and $\tilde{X} = X - \mu_X$:
\[
\tilde{Y} = \tilde{X}'\beta + e.
\]

The FOC for the demeaned problem gives:
\begin{align*}
\E[\tilde{X}\tilde{Y}] &= \E[\tilde{X}\tilde{X}']\beta\\
\Cov(X,Y) &= \Var(X) \cdot \beta.
\end{align*}

Here:
\[
\Var(X) = \E[(X - \mu_X)(X - \mu_X)'] =
\begin{pmatrix}
\Var(X_1) & \Cov(X_1,X_2) & \cdots \\
\Cov(X_2,X_1) & \Var(X_2) & \cdots \\
\vdots & \vdots & \ddots
\end{pmatrix}
\]
and
\[
\Cov(X,Y) =
\begin{pmatrix}
\Cov(X_1, Y) \\
\Cov(X_2, Y) \\
\vdots
\end{pmatrix}.
\]

Solving:
\begin{equation}
\boxed{\beta = \Var(X)^{-1}\Cov(X,Y).}
\end{equation}

\textbf{Connection to the second-moment formula.}  The full formula $\beta = \E[XX']^{-1}\E[XY]$ with the constant included is \emph{algebraically equivalent} to computing the slope via centered moments and then recovering the intercept. Demeaning ``absorbs'' the intercept.


%======================================================================
\section{Partitioned Regression}
\label{sec:partitioned}
%======================================================================

\begin{tcolorbox}[colback=blue!5, colframe=blue!50, title=Slide Reference]
``Partitions and Regression Sub-Vectors,'' ``Partitioned Matrix Inversion and Regression,'' ``Interpreting Multivariate Regression,'' ``Demonstration: Iterated Projection''
\end{tcolorbox}

\subsection{Setup}

Partition $X$ into two sub-vectors:
\[
X = \begin{pmatrix} X_1 \\ X_2 \end{pmatrix}, \quad
\beta = \begin{pmatrix} \beta_1 \\ \beta_2 \end{pmatrix},
\]
where $X_1$ is $k_1 \times 1$ and $X_2$ is $k_2 \times 1$ (with $k_1 + k_2 = k$).

For concreteness, suppose $k_1 = 1$ (a single variable of interest) and $k_2 = 2$ (two controls plus the constant):
\[
X_1 = (X_1), \quad X_2 = \begin{pmatrix} X_2 \\ X_3 \\ 1 \end{pmatrix}.
\]

The linear projection is:
\[
Y = X_1'\beta_1 + X_2'\beta_2 + e, \quad \E[Xe] = 0.
\]

\subsection{Partitioning $\bfQ_{XX}$ and $\bfQ_{XY}$}

The normal equation $\bfQ_{XY} = \bfQ_{XX}\beta$ becomes:
\[
\begin{pmatrix} \bfQ_{1Y} \\ \bfQ_{2Y} \end{pmatrix}
=
\begin{pmatrix} \bfQ_{11} & \bfQ_{12} \\ \bfQ_{21} & \bfQ_{22} \end{pmatrix}
\begin{pmatrix} \beta_1 \\ \beta_2 \end{pmatrix},
\]
where:
\begin{align*}
\bfQ_{11} &= \E[X_1 X_1'] \quad (k_1 \times k_1), &
\bfQ_{12} &= \E[X_1 X_2'] \quad (k_1 \times k_2), \\
\bfQ_{21} &= \E[X_2 X_1'] \quad (k_2 \times k_1), &
\bfQ_{22} &= \E[X_2 X_2'] \quad (k_2 \times k_2).
\end{align*}

This is two equations:
\begin{align}
\bfQ_{1Y} &= \bfQ_{11}\beta_1 + \bfQ_{12}\beta_2, \label{eq:part1}\\
\bfQ_{2Y} &= \bfQ_{21}\beta_1 + \bfQ_{22}\beta_2. \label{eq:part2}
\end{align}

\subsection{Solving for $\beta_1$ via elimination}

\begin{step}[Solve \eqref{eq:part2} for $\beta_2$]
\[
\beta_2 = \bfQ_{22}^{-1}(\bfQ_{2Y} - \bfQ_{21}\beta_1).
\]
\end{step}

\begin{step}[Substitute into \eqref{eq:part1}]
\begin{align*}
\bfQ_{1Y} &= \bfQ_{11}\beta_1 + \bfQ_{12}\bfQ_{22}^{-1}(\bfQ_{2Y} - \bfQ_{21}\beta_1)\\
\bfQ_{1Y} - \bfQ_{12}\bfQ_{22}^{-1}\bfQ_{2Y} &= (\bfQ_{11} - \bfQ_{12}\bfQ_{22}^{-1}\bfQ_{21})\beta_1.
\end{align*}
\end{step}

\begin{step}[Identify the ``residualized'' quantities]
Define:
\begin{align*}
\bfQ_{11\cdot 2} &\equiv \bfQ_{11} - \bfQ_{12}\bfQ_{22}^{-1}\bfQ_{21}, \\
\bfQ_{1Y\cdot 2} &\equiv \bfQ_{1Y} - \bfQ_{12}\bfQ_{22}^{-1}\bfQ_{2Y}.
\end{align*}
Then:
\begin{equation}
\boxed{\beta_1 = \bfQ_{11\cdot 2}^{-1}\bfQ_{1Y\cdot 2}.} \label{eq:beta1part}
\end{equation}
By symmetry: $\beta_2 = \bfQ_{22\cdot 1}^{-1}\bfQ_{2Y\cdot 1}$.
\end{step}

\subsection{What $\bfQ_{11\cdot 2}$ and $\bfQ_{1Y\cdot 2}$ mean: the iterated projection}

\begin{tcolorbox}[colback=blue!5, colframe=blue!50, title=Slide Reference]
``Demonstration: Iterated Projection''
\end{tcolorbox}

\paragraph{Step A: Regress $X_1$ on $X_2$.}
\[
X_1 = X_2'\gamma_2 + u_1, \quad \E[X_2 u_1] = 0.
\]
The coefficient is $\gamma_2 = \bfQ_{22}^{-1}\bfQ_{21}$ and the residual $u_1 = X_1 - X_2'\gamma_2$ is the part of $X_1$ that $X_2$ cannot predict.

\paragraph{Step B: Compute $\E[u_1^2]$.}
\begin{align*}
\E[u_1^2] &= \E[(X_1 - X_2'\gamma_2)^2]\\
&= \E[X_1^2] - 2\gamma_2'\E[X_2 X_1] + \gamma_2'\E[X_2 X_2']\gamma_2\\
&= \bfQ_{11} - 2\bfQ_{21}'\bfQ_{22}^{-1}\bfQ_{21} + \bfQ_{21}'\bfQ_{22}^{-1}\bfQ_{22}\bfQ_{22}^{-1}\bfQ_{21}\\
&= \bfQ_{11} - \bfQ_{12}\bfQ_{22}^{-1}\bfQ_{21} = \bfQ_{11\cdot 2}.
\end{align*}
So $\bfQ_{11\cdot 2}$ is the variance of $X_1$ after removing what $X_2$ predicts.

\paragraph{Step C: Compute $\E[u_1 Y]$.}
\begin{align*}
\E[u_1 Y] &= \E[(X_1 - X_2'\gamma_2)Y]\\
&= \E[X_1 Y] - \gamma_2'\E[X_2 Y]\\
&= \bfQ_{1Y} - \bfQ_{12}\bfQ_{22}^{-1}\bfQ_{2Y} = \bfQ_{1Y\cdot 2}.
\end{align*}
So $\bfQ_{1Y\cdot 2}$ is the covariance of the residualized $X_1$ with $Y$.

\paragraph{Conclusion:}
\begin{equation}
\beta_1 = \bfQ_{11\cdot 2}^{-1}\bfQ_{1Y\cdot 2} = \frac{\E[u_1 Y]}{\E[u_1^2]}.
\end{equation}

When $X_1$ is a single variable (scalar), this is just $\Cov(u_1, Y)/\Var(u_1)$: the bivariate slope of $Y$ on the part of $X_1$ that is unique---not explained by $X_2$.

\textbf{This is the Frisch--Waugh--Lovell intuition:} The coefficient on $X_1$ in a multivariate regression is identical to the coefficient from a simple regression of $Y$ on the residuals from regressing $X_1$ on all other regressors.


%======================================================================
\section{Omitted Variable Bias}
\label{sec:ovb}
%======================================================================

\begin{tcolorbox}[colback=blue!5, colframe=blue!50, title=Slide Reference]
``Omitted Variable Bias''
\end{tcolorbox}

\subsection{Setup}

\textbf{Long regression} (the correctly specified model):
\[
Y = X_1'\beta_1 + X_2'\beta_2 + e, \quad \E[Xe] = 0.
\]

\textbf{Short regression} (omitting $X_2$):
\[
Y = X_1'\gamma_1 + u, \quad \E[X_1 u] = 0.
\]

The short-regression coefficient is:
\[
\gamma_1 = (\E[X_1 X_1'])^{-1}\E[X_1 Y].
\]

\subsection{Deriving the bias}

Substitute $Y = X_1'\beta_1 + X_2'\beta_2 + e$ into the formula for $\gamma_1$:

\begin{align*}
\gamma_1 &= (\E[X_1 X_1'])^{-1}\E[X_1(X_1'\beta_1 + X_2'\beta_2 + e)]\\
&= (\E[X_1 X_1'])^{-1}\bigl[\E[X_1 X_1']\beta_1 + \E[X_1 X_2']\beta_2 + \E[X_1 e]\bigr].
\end{align*}

\textbf{Term by term:}

\paragraph{First term:}
$(\E[X_1 X_1'])^{-1}\E[X_1 X_1']\beta_1 = \mathbf{I}_{k_1}\beta_1 = \beta_1.$

\paragraph{Second term:}
$(\E[X_1 X_1'])^{-1}\E[X_1 X_2'] = \Gamma_{12}$, a $k_1 \times k_2$ matrix.

\textbf{What is $\Gamma_{12}$?}  It is the matrix of coefficients from regressing each variable in $X_2$ on $X_1$. The $(i,j)$ entry tells you how much the $j$-th element of $X_2$ is associated with the $i$-th element of $X_1$.

\paragraph{Third term:}
$(\E[X_1 X_1'])^{-1}\E[X_1 e] = (\E[X_1 X_1'])^{-1}\cdot \mathbf{0} = \mathbf{0}$, because $\E[X_1 e] = 0$ from the long regression (the first $k_1$ rows of $\E[Xe] = \mathbf{0}$).

\paragraph{Result:}
\begin{equation}
\boxed{\gamma_1 = \beta_1 + \Gamma_{12}\beta_2.}
\end{equation}

The \textbf{omitted variable bias} is $\Gamma_{12}\beta_2$.

\subsection{Expanded for the scalar case}

If $X_1$ and $X_2$ are both scalars:
\[
\Gamma_{12} = \frac{\E[X_1 X_2]}{\E[X_1^2]} = \frac{\Cov(X_1, X_2)}{\Var(X_1)} \quad \text{(when demeaned)},
\]
and the bias is:
\[
\text{Bias} = \frac{\Cov(X_1, X_2)}{\Var(X_1)} \cdot \beta_2.
\]

\begin{itemize}
\item If $X_1$ and $X_2$ are positively correlated ($\Gamma_{12} > 0$) and $X_2$ has a positive effect on $Y$ ($\beta_2 > 0$), then $\gamma_1 > \beta_1$: the short regression \textbf{overstates} the effect.
\item If the correlation and the effect have opposite signs, the bias goes the other direction.
\item If $X_1$ and $X_2$ are uncorrelated ($\Gamma_{12} = 0$), there is no bias: $\gamma_1 = \beta_1$.
\end{itemize}


%======================================================================
\section{CEF Algebra: Variance Decomposition and Theorem Proofs}
\label{sec:cef}
%======================================================================

\begin{tcolorbox}[colback=blue!5, colframe=blue!50, title=Slide Reference]
``Covariance of CEF and Its Error,'' ``Theorem 2.7: CEF as Best Predictor''
\end{tcolorbox}

The CEF derivations are primarily probability rather than linear algebra, but they use algebraic expansions that parallel the BLP derivations.

\subsection{Covariance of $m(X)$ and $e$}

We show $\Cov(e, m(X)) = 0$ where $e = Y - m(X)$ and $m(X) = \E[Y|X]$.

\begin{step}[Compute $\E[e \cdot m(X)]$]
Use the law of iterated expectations, conditioning on $X$:
\begin{align*}
\E[e \cdot m(X)] &= \E\!\big[\E[e \cdot m(X) | X]\big]\\
&= \E\!\big[m(X) \cdot \E[e|X]\big] \quad \text{($m(X)$ is known given $X$, so it factors out)}\\
&= \E[m(X) \cdot 0] = 0.
\end{align*}
\end{step}

\begin{step}[Compute the covariance]
\begin{align*}
\Cov(e, m(X)) &= \E[e \cdot m(X)] - \E[e]\E[m(X)]\\
&= 0 - 0 \cdot \E[m(X)] = 0.
\end{align*}
(We used $\E[e] = 0$ from the LIE.)
\end{step}

\begin{step}[Variance decomposition]
Since $Y = m(X) + e$ and $\Cov(m(X), e) = 0$:
\[
\Var(Y) = \Var(m(X)) + \Var(e) + 2\Cov(m(X), e) = \Var(m(X)) + \Var(e).
\]
\end{step}

\subsection{Theorem 2.7: CEF is the best predictor}

We show $\E[(Y - g(X))^2] \geq \E[(Y - m(X))^2]$ for any function $g$.

\begin{step}[Decompose using $Y = m(X) + e$]
\[
Y - g(X) = e + (m(X) - g(X)).
\]
\end{step}

\begin{step}[Expand the square]
\begin{align*}
\E[(Y - g(X))^2] &= \E\!\big[(e + (m(X) - g(X)))^2\big]\\
&= \E[e^2] + 2\E[e(m(X) - g(X))] + \E[(m(X) - g(X))^2].
\end{align*}
\end{step}

\begin{step}[The cross term vanishes]
Since $m(X) - g(X)$ is a function of $X$, Theorem 2.4.4 gives $\E[e \cdot h(X)] = 0$ for any measurable $h(X)$. So $\E[e(m(X) - g(X))] = 0$.
\end{step}

\begin{step}[Conclude]
\[
\E[(Y - g(X))^2] = \underbrace{\E[e^2]}_{\E[(Y-m(X))^2]} + \underbrace{\E[(m(X) - g(X))^2]}_{\geq\, 0} \;\geq\; \E[(Y - m(X))^2].
\]
Equality holds if and only if $g(X) = m(X)$ almost surely.
\end{step}


%======================================================================
\section{Partitioned Matrix Inverse Formula}
\label{sec:blockinverse}
%======================================================================

\begin{tcolorbox}[colback=blue!5, colframe=blue!50, title=Slide Reference]
``Partitioned Matrix Inversion and Regression''
\end{tcolorbox}

The slide states the block inverse of $\bfQ_{XX}$:
\[
\bfQ_{XX}^{-1} =
\begin{pmatrix}
\bfQ_{11\cdot 2}^{-1} & -\bfQ_{11\cdot 2}^{-1}\bfQ_{12}\bfQ_{22}^{-1} \\
-\bfQ_{22\cdot 1}^{-1}\bfQ_{21}\bfQ_{11}^{-1} & \bfQ_{22\cdot 1}^{-1}
\end{pmatrix},
\]
where $\bfQ_{11\cdot 2} = \bfQ_{11} - \bfQ_{12}\bfQ_{22}^{-1}\bfQ_{21}$ and $\bfQ_{22\cdot 1} = \bfQ_{22} - \bfQ_{21}\bfQ_{11}^{-1}\bfQ_{12}$.

\subsection{Verification}

We verify by checking $\bfQ_{XX} \cdot \bfQ_{XX}^{-1} = \mathbf{I}$. Denote the proposed inverse as $\mathbf{M}$:
\[
\begin{pmatrix} \bfQ_{11} & \bfQ_{12} \\ \bfQ_{21} & \bfQ_{22} \end{pmatrix}
\begin{pmatrix} \mathbf{M}_{11} & \mathbf{M}_{12} \\ \mathbf{M}_{21} & \mathbf{M}_{22} \end{pmatrix}
= \begin{pmatrix} \mathbf{I}_{k_1} & \mathbf{0} \\ \mathbf{0} & \mathbf{I}_{k_2} \end{pmatrix}.
\]

\paragraph{Upper-left block:}
\begin{align*}
\bfQ_{11}\bfQ_{11\cdot 2}^{-1} + \bfQ_{12}(-\bfQ_{22\cdot 1}^{-1}\bfQ_{21}\bfQ_{11}^{-1}).
\end{align*}
This requires extended algebra involving the Schur complement identity. The key fact is:
\[
\bfQ_{11\cdot 2}^{-1} = (\bfQ_{11} - \bfQ_{12}\bfQ_{22}^{-1}\bfQ_{21})^{-1},
\]
which is the Schur complement of $\bfQ_{22}$ in the block matrix. The formula is a standard result from block matrix algebra.

\paragraph{Upper-right block (checking it equals $\mathbf{0}$):}
\begin{align*}
&\bfQ_{11}(-\bfQ_{11\cdot 2}^{-1}\bfQ_{12}\bfQ_{22}^{-1}) + \bfQ_{12}\bfQ_{22\cdot 1}^{-1}\\
&= -\bfQ_{11}\bfQ_{11\cdot 2}^{-1}\bfQ_{12}\bfQ_{22}^{-1} + \bfQ_{12}\bfQ_{22\cdot 1}^{-1}.
\end{align*}
Using the Schur complement identity $\bfQ_{11}\bfQ_{11\cdot 2}^{-1} = \mathbf{I} + \bfQ_{12}\bfQ_{22}^{-1}\bfQ_{21}\bfQ_{11\cdot 2}^{-1}$ and further simplification via the Woodbury identity confirms this equals $\mathbf{0}$.

\subsection{Why this matters for regression}

Multiplying the block inverse by $\bfQ_{XY}$:
\begin{align*}
\begin{pmatrix} \beta_1 \\ \beta_2 \end{pmatrix}
&= \bfQ_{XX}^{-1}\bfQ_{XY}
= \begin{pmatrix}
\bfQ_{11\cdot 2}^{-1} & -\bfQ_{11\cdot 2}^{-1}\bfQ_{12}\bfQ_{22}^{-1} \\
-\bfQ_{22\cdot 1}^{-1}\bfQ_{21}\bfQ_{11}^{-1} & \bfQ_{22\cdot 1}^{-1}
\end{pmatrix}
\begin{pmatrix} \bfQ_{1Y} \\ \bfQ_{2Y} \end{pmatrix}.
\end{align*}

For $\beta_1$:
\begin{align*}
\beta_1 &= \bfQ_{11\cdot 2}^{-1}\bfQ_{1Y} - \bfQ_{11\cdot 2}^{-1}\bfQ_{12}\bfQ_{22}^{-1}\bfQ_{2Y}\\
&= \bfQ_{11\cdot 2}^{-1}\bigl(\bfQ_{1Y} - \bfQ_{12}\bfQ_{22}^{-1}\bfQ_{2Y}\bigr)\\
&= \bfQ_{11\cdot 2}^{-1}\bfQ_{1Y\cdot 2},
\end{align*}
confirming equation~\eqref{eq:beta1part}.


%======================================================================
\section{Summary of Key Formulas}
\label{sec:summary}
%======================================================================

\begin{center}
\renewcommand{\arraystretch}{1.8}
\begin{tabular}{p{4cm} l l}
\toprule
\textbf{Quantity} & \textbf{Bivariate} & \textbf{Multivariate} \\
\midrule
BLP slope & $\dfrac{\Cov(X,Y)}{\Var(X)}$ & $\E[XX']^{-1}\E[XY]$ \\[6pt]
BLP intercept & $\E[Y] - \beta\E[X]$ & (absorbed into $X$ when $1 \in X$) \\[6pt]
Orthogonality & $\E[Xe] = 0$ & $\E[Xe] = \mathbf{0}_{k\times 1}$ \\[6pt]
Error variance & $\Var(Y) - \beta^2\Var(X)$ & $Q_{YY} - \bfQ_{YX}\bfQ_{XX}^{-1}\bfQ_{XY}$ \\[6pt]
Sub-vector coeff. & --- & $\bfQ_{11\cdot 2}^{-1}\bfQ_{1Y\cdot 2}$ \\[6pt]
OVB & $\dfrac{\Cov(X_1,X_2)}{\Var(X_1)}\beta_2$ & $\Gamma_{12}\beta_2$ \\
\bottomrule
\end{tabular}
\end{center}

\bigskip
In each case, the multivariate formula reduces to the bivariate formula when $k=1$ (plus a constant). The matrix algebra is the same operation---projecting and residualizing---applied to vectors instead of scalars.

\end{document}
